\section{Introduction to Groups}

\subsection{Basic Axioms and Examples}

\begin{groupdef}
    \defn{Binary Operation} On a set $G$, it is the function $\star : G \times G \to G$. We write $a \star b$ for $\star(a, b)$, where $a, b \in G$.

    \defn{Associative} A binary operation $\star$ on a set $G$ such that $a \star (b \star c) = (a \star b) \star c$ for all $a, b, c \in G$.

    \defn{Commutative} A binary operation $\star$ on a set $G$ such that $a \star b = b \star a$ for all $a, b \in G$.
    If $a \star b = b \star a$ for \emph{some} $a, b \in G$, we say that $a$ and $b$ \term{commute}.
\end{groupdef}

\begin{example}[Binary Operations]
    \begin{enumerate}
        \item $+$ and $\cdot$ are associative and commutative binary operations on $\zbb$, $\qbb$, $\rbb$, and $\cbb$.
        \item $-$ is not a binary operation on $\zp$ since $a - b \notin \zp$ when $a < b$, i.e., $- : \zp \times \zp \to \zp$ is not well-defined.
        \item The operation of taking cross products in $\rbb^3$ is a non-associative and non-commutative binary operation on $\rbb^3$.
    \end{enumerate}
\end{example}

\defn{Closure} Given a binary operation $\star$ on a set $G$ with subset $H \subseteq G$, we say that $H$ is \term{closed under $\star$} if $a \star b \in H$ for all $a, b \in H$.

\defn{Group} An ordered pair $(G, \star)$ for a set $G$ and a binary operation $\star$ on $G$ such that
\begin{enumerate}
    \item $\star$ is associative, i.e., $a \star (b \star c) = (a \star b) \star c$ for all $a, b, c \in G$,
    \item there exists an \term{identity element} $e \in G$ such that $e \star a = a \star e = a$ for all $a \in G$, and
    \item for each $a \in G$, there exists $a\inv \in G$, called the \term{inverse} of $a$, such that $a \star a\inv = a\inv \star a = e$.
\end{enumerate}

\defn{Abelian} A group $(G, \star)$ is \term{Abelian} if $\star$ is commutative.

\defn{Finite Group} A group $(G, \star)$ is \term{finite} if $G$ is a finite set.

\begin{example}[Groups]
    \begin{enumerate}
        \item $(G, +)$ is an Abelian group for $G = \zbb$, $\qbb$, $\rbb$, and $\cbb$ with $e = 0$ and $a\inv = -a$ for all $a \in G$.
        \item The axioms of a vector space $V$ imply that $(V, +)$ is an Abelian group for all $n \in \zp$.
    \end{enumerate}
\end{example}

\defn{Direct Product} Given groups $(G, \star)$ and $(H, \diamond)$, their \term{direct product} is the group $(G \times H, \ast)$ where
\[
    G \times H = \set{(g, h) \mid g \in G, h \in H}
\]
and $\ast$ is defined for all $(g_1, h_1), (g_2, h_2) \in G \times H$ as
\[
    (g_1, h_1) \ast (g_2, h_2) = (g_1 \star g_2, h_1 \diamond h_2).
\]

\begin{proposition}[Group Properties]
    Let $(G, \star)$ be a group.
    \begin{enumerate}
        \item The identity element $e \in G$ is unique.
        \item For each $g \in G$, the inverse $g\inv \in G$ is unique.
        \item $(g\inv)\inv = g$ for all $g \in G$.
        \item For all $g, h \in G$, $(g \star h)\inv = h\inv \star g\inv$.
        \item For any $g_1, \ldots, g_n \in G$, the value of $g_1 \star g_2 \star \cdots \star g_n$ is independent of how the product is parenthesized, also known as the generalized associative law.
    \end{enumerate}
\end{proposition}

\begin{proofnum}
    \item Let $e_1, e_2 \in G$ be identity elements. Then
    \[
        e_1 = e_1 \star e_2 = e_2.
    \]
    \item Let $g \in G$ and let $g\inv_1, g\inv_2 \in G$ be inverses of $g$. Then
    \[
        g\inv_1 = g\inv_1 \star e = g\inv_1 \star (g \star g\inv_2) = (g\inv_1 \star g) \star g\inv_2 = e \star g\inv_2 = g\inv_2.
    \]
    \item Note that $g\inv(g\inv)\inv = e = g\inv g$. By the uniqueness of inverses, we have $(g\inv)\inv = g$.
    \item Let $k = (g \star h)\inv$. Then
    \[
        (g \star h) \star k = e.
    \]
    Left multiply both sides by $g\inv$ and use the associativity of $\star$ to obtain
    \begin{align*}
        g\inv \star ((g \star h) \star k) &= g\inv \star e \\
        (g\inv \star g) \star (h \star k) &= g\inv \\
        e \star (h \star k) &= g\inv \\
        h \star k &= g\inv.
    \end{align*}
    Left multiply both sides by $h\inv$ and simplify similarly to obtain
    \begin{align*}
        h\inv \star (h \star k) &= h\inv \star g\inv \\
        (h\inv \star h) \star k &= h\inv \star g\inv \\
        e \star k &= h\inv \star g\inv \\
        k &= h\inv \star g\inv.
    \end{align*}
    \item We proceed by induction on $n$.
    The base case $n = 1$ is trivial.
    For cases $n = 2$ and $n = 3$, the result follows from the associativity of $\star$.
    Suppose for all $k < n$ that the bracketing of a product of the $k$ elements $h_1 \star \cdots \star h_k$ can be reduced to the form
    \[
        h_1 \star (h_2 \star (\cdots \star (h_{k-1} \star h_k)) \cdots).
    \]
    
    For notation, write the product above as $R(h_1, \ldots, h_k)$.
    Consider a product of $n$ elements $g_1 \star \cdots \star g_n$ with some arbitrary bracketing.
    Then we can write the product as
    \[
        \underbrace{(g_1 \star \cdots \star g_i)}_{i \text{ factors}} \star \underbrace{(g_{i + 1} \star \cdots \star g_n)}_{n-i \text{ factors}}
    \]
    for some $1 \leq i < n$.
    By the inductive hypothesis, we may write the $i$ factors as $R(g_1, \ldots, g_i)$ and the $n-i$ factors as $R(g_{i + 1}, \ldots, g_n)$.
    Thus, we have
    \[
        g_1 \star \cdots \star g_n = R(g_1, \ldots, g_i) \star R(g_{i + 1}, \ldots, g_n).
    \]

    We first show that
    \[
        R(g_1, \ldots, g_i) \star x = R(g_1, \ldots, g_i, x) \text{ for any } x \in G.
    \]
    We proceed by induction on $i$.
    The base case $i = 1$ is trivial.
    For $i = 2$, we have 
    \[
        R(g_1, g_2) \star x = (g_1 \star g_2) \star x = g_1 \star (g_2 \star x) = R(g_1, g_2, x) \text{ by the associativity of } \star.
    \]
    Suppose for all $j < i$ that $R(g_1, \ldots, g_j) \star x = R(g_1, \ldots, g_j, x)$ for any $x \in G$.
    Noting that $R(g_1, \ldots, g_i) = g_1 \star R(g_2, \ldots, g_i)$, we have
    \begin{align*}
        R(g_1, \ldots, g_i) \star x &= (g_1 \star R(g_2, \ldots, g_i)) \star x \\
        &= g_1 \star (R(g_2, \ldots, g_i) \star x) \\
        &= g_1 \star R(g_2, \ldots, g_i, x) \\
        &= R(g_1, \ldots, g_i, x),
    \end{align*}
    where the second equality follows from the associativity of $\star$ and the third equality follows from the inductive hypothesis.

    Using the above claim and putting $x = R(g_{i + 1}, \ldots, g_n)$, we have
    \begin{align*}
        g_1 \star \cdots \star g_n &= R(g_1, \ldots, g_i) \star R(g_{i + 1}, \ldots, g_n) \\
        &= R(g_1, \ldots, g_i, R(g_{i + 1}, \ldots, g_n)) \\
        &= R(g_1, \ldots, g_n),
    \end{align*}
    where the last equality follows from the definition of $R$.
    Thus, by induction, the value of $g_1 \star \cdots \star g_n$ is independent of how the product is parenthesized.
\end{proofnum}

\begin{note}
    We may drop the $\star$ symbol and write $gh$ for $g \star h$ when the binary operation is clear from context.
    Moreover, we may write $g^n$ for the product of $n$ copies of $g \in G$, where $g^0 = e$ and $g^{-n} = (g\inv)^n$ for all $n \in \zp$.
    Additionally, we may write $1$ for the identity element $e \in G$.
\end{note}

\begin{proposition}[Left and Right Cancellation Laws]
    Let $G$ be a group with $g, h \in G$.
    Then the equations $gs = h$ and $tg = h$ have unique solutions $s, t \in G$, respectively.
    In particular,
    \begin{enumerate}
        \item if $gx = gy$, then $x = y$, and
        \item if $xh = yh$, then $x = y$.
    \end{enumerate}
\end{proposition}

\begin{proof}
    The equation $gs = h$ can be solved by left multiplying both sides by $g\inv$ to obtain $s = g\inv h$, which is unique by the uniqueness of inverses.
    Similarly, $tg = h$ can be solved to obtain the unique solution $t = hg\inv$.

    The first cancellation law follows from left multiplying both sides of $gx = gy$ by $g\inv$ to obtain $x = y$, and the second cancellation law follows from right multiplying both sides of $xh = yh$ by $h\inv$ to obtain $x = y$.
\end{proof}

\defn{Order} The smallest $n \in \zp$ for a group element $g \in G$ such that $g^n = 1$, denoted as $|g| = n$.
If no such $n$ exists, we say that $g$ has \term{infinite order} and write $|g| = \infty$.

\defn{Multiplication/Group Table} Put $G = \set{g_1, g_2, \ldots, g_n}$ with $g_1 = 1$.
Then the \term{multiplication table} or \term{group table} of $G$ is the $n \times n$ matrix whose $(i, j)$th entry is $g_i g_j$ for all $1 \leq i, j \leq n$.